\documentclass[11pt]{article}
\usepackage[margin=1in]{geometry}
\usepackage[authoryear,sort]{natbib}
\bibliographystyle{apalike}
\usepackage{amsmath,amssymb,amsthm}
\newcommand{\Pcal}{\mathcal{P}}
\newcommand{\Hcal}{\mathcal{H}}
\newcommand{\Fcal}{\mathcal{F}}
\newcommand{\EE}{\mathbb{E}}
\newcommand{\PP}{\mathbb{P}}
\newcommand{\DRL}{\mathrm{DRL}}
\newcommand{\DRU}{\mathrm{DRU}}
\newcommand{\LIL}{\mathrm{LIL}}

\newtheorem{theorem}{Theorem}
\newtheorem{lemma}{Lemma}
\newtheorem{assumption}{Assumption}

\begin{document}

\section{Setup and notation}

Let $\Pi=\{\pi_1,\dots,\pi_m\}$ denote a finite class of policies, where $m:=|\Pi|$. 
We observe contextual bandit data $(X_t,A_t,R_t)_{t\ge1}$ generated sequentially by a predictable sequence of logging policies $(h_t)_{t\ge1}$. 
Formally, conditional on the past history $\mathcal H_{t-1}$ and the current context $X_t$, the action satisfies
\[
A_t \mid (X_t,\mathcal H_{t-1}) \sim h_t(\cdot \mid X_t),
\]
where $\mathcal H_{t-1}$ denotes the full observed history up to time $t-1$. 
We assume the rewards satisfy $R_t\in[0,1]$ almost surely.

\noindent For a target policy $\pi\in\Pi$, define the importance weight
\[
w_t(\pi):=\frac{\pi(A_t\mid X_t)}{h_t(A_t\mid X_t)}.
\]

\noindent We work in the time-invariant policy value setting: for each $\pi\in\Pi$, let
\[
\nu_t(\pi):=\mathbb{E}_\pi[R_t | \mathcal H_{t-1}],
\]
and assume that there exists a constant value $\nu(\pi)\in[0,1]$ such that $\nu_t(\pi)=\nu(\pi)$ for all $t\ge1$.
Let
\[
\pi^\star \in \arg\max_{\pi\in\Pi} \nu(\pi)
\]
denote an optimal policy.

\paragraph{Doubly robust pseudo-outcomes.}
Let $\widehat r_t(x,a)\in[0,1]$ denote a predictable reward predictor constructed from the history $\mathcal H_{t-1}$. Define the lower and upper doubly robust pseudo-outcomes
\begin{align}
\phi_t^{\mathrm{DRL}}(\pi)
&:= w_t(\pi)\!\left(R_t-\Big[\widehat r_t(X_t,A_t)\wedge \tfrac{k_t}{w_t(\pi)}\Big]\right) \notag\\
&\quad + \mathbb{E}_{a\sim\pi(\cdot\mid X_t)}
\!\left(\widehat r_t(X_t,a)\wedge \tfrac{k_t}{w_t(\pi)}\right), \\
\phi_t^{\mathrm{DRU}}(\pi)
&:= w_t(\pi)\!\left((1-R_t)-\Big[(1-\widehat r_t(X_t,A_t))\wedge \tfrac{k_t}{w_t(\pi)}\Big]\right) \notag\\
&\quad + \mathbb{E}_{a\sim\pi(\cdot\mid X_t)}
\!\left((1-\widehat r_t(X_t,a))\wedge \tfrac{k_t}{w_t(\pi)}\right).
\end{align}

\noindent As shown by \citet{waudby-smith_anytime-valid_2024}, these pseudo-outcomes satisfy
\[
\EE\!\left[\phi_t^{\DRL}(\pi)\mid \mathcal H_{t-1}\right]=\nu(\pi),\qquad
\EE\!\left[\phi_t^{\DRU}(\pi)\mid \mathcal H_{t-1}\right]=1-\nu(\pi),
\]
and are bounded below by $-k_t$.

\paragraph{CS construction.}
Let $k_1=k_2=\cdots = k \ge 0$ be a fixed constant and let $\xi_0 \in [0,(1+k)^{-1}]$. For $t \ge 1$, define the scaled doubly robust pseudo-outcomes
\[
\xi_t := \frac{\phi_t}{1+k},
\]
where $\phi_t := \phi_t^{\DRL}$ or $\phi_t^{\DRU}$, depending on which bound is being constructed. Then define
\[
\hat{\xi}_0 := \xi_0 \ ; \ \hat{\xi}_t := \left(\frac{1}{t}\sum_{i=1}^t \xi_i\right)\wedge \frac{1}{1+k},
\]
the variance process
\[
V_t := \sum_{i=1}^t (\xi_i - \hat{\xi}_{i-1})^2,
\]
and the stabilized version
\[
\bar V_t := V_t \vee 1.
\]

\noindent Let $\alpha \in (0,1)$ denote a significance level and define
\[
\ell_t(\alpha) := 2\log(\log \bar V_t + 1) + \log\left(\frac{1.65}{\alpha}\right).
\]

\noindent Using these quantities, define the lower CS boundary for a particular policy $\pi$ and significance level $\alpha$:
\[
L_t^{\LIL}(\pi; \alpha)
:= (k+1)\left(
\frac{1}{t}\sum_{i=1}^t \xi_i
-
\frac{\sqrt{2.13\,\ell_t(\alpha)\,\bar V_t + 1.76\,\ell_t(\alpha)^2}}{t}
-
\frac{1.33\,\ell_t(\alpha)^2}{t}
\right)\lor 0.
\]

\noindent $U_t^{\LIL}(\pi;\alpha)$ denotes the corresponding upper confidence-sequence bound obtained analogously by the mirrored construction using $\phi_t^{\DRU}$. From here on, let $\left[L_t^{\LIL}(\pi;\alpha), \ U_t^{\LIL}(\pi;\alpha)\right]$ denote a valid two-sided $(1-\alpha)$ CS constructed by a union bound over the lower and upper one-sided CSs, $L_t^{\LIL}(\pi;\alpha/2)$ and $U_t^{\LIL}(\pi;\alpha/2)$.

\section{Optimal policy set}

With a fixed, overall significance level $\alpha$, set the policy-level significance level to $\frac{\alpha}{m}$ for each $\pi\in\Pi$. It follows by a union bound over $\pi \in \Pi$ that 
\begin{equation}
\PP\left(\nu(\pi) \in \left[L_t^\LIL\left(\pi; \frac{\alpha}{m}\right), U_t^\LIL\left(\pi; \frac{\alpha}{m}\right)\right]: \forall t \ge 1, \forall \pi \in \Pi \right) \ge 1 - \alpha.
\label{eq:simultaneous_coverage}
\end{equation}

\begin{lemma}[Anytime-valid optimal policy set]
Define the set
\[
S_t:=\left\{\pi\in\Pi:\ U_t^{\LIL}\left(\pi; \frac{\alpha}{m}\right)\ \ge\ \max_{\pi'\in\Pi} L_t^{\LIL}\left(\pi'; \frac{\alpha}{m}\right)\right\}
\]
and the set of optimal policies
\[
\Pi^* := \left\{\pi \in \Pi : \nu(\pi) = \underset{\pi'\in\Pi}{\max}\ \nu(\pi') \right\}.
\]
Then,
\[
\PP\left(\Pi^* \subseteq S_t : \forall t \ge 1\right) \ge 1 - \alpha.
\]
\end{lemma}

\begin{proof}
For each $t \ge 1$, define
\[
E(t) := \left\{\nu(\pi) \in \left[L_t^\LIL\left(\pi; \frac{\alpha}{m}\right), U_t^\LIL\left(\pi; \frac{\alpha}{m}\right)\right] : \forall \pi \in \Pi \right\},
\]
and let
\[
E := \bigcap_{t\ge1} E(t).
\]
As established above, $\PP(E) \ge 1-\alpha$. Now fix any $\pi^\dagger \in \Pi^*$ and any $t \ge 1$. If $E(t)$ holds true, we have that
\[
L_t^\LIL\left(\pi'; \frac{\alpha}{m}\right) \le \nu(\pi') \le \nu(\pi^\dagger)
\qquad \text{for all } \pi' \in \Pi,
\]
since $\pi^\dagger$ is optimal. Therefore,
\[
\max_{\pi'\in\Pi} L_t^\LIL\left(\pi'; \frac{\alpha}{m}\right) \le \nu(\pi^\dagger) \le U_t^\LIL\left(\pi^\dagger; \frac{\alpha}{m}\right).
\]
Combining these gives
\[
\max_{\pi'\in\Pi} L_t^\LIL\left(\pi'; \frac{\alpha}{m}\right)
\le
U_t^\LIL\left(\pi^\dagger; \frac{\alpha}{m}\right),
\]
so $\pi^\dagger \in S_t$. Since $\pi^\dagger \in \Pi^*$ was arbitrary, it follows that when $E(t)$ holds, $\Pi^* \subseteq S_t$. Thus,
\[
E \implies \Pi^* \subseteq S_t \ \ \forall t\ge1.
\]
Finally,
\[
\PP\left(\Pi^* \subseteq S_t : \forall t \ge 1\right)
\ge
\PP(E)
\ge
1-\alpha.
\]
\end{proof}

\section{CS width asymptotic rates}

By Proposition 3 and the accompanying discussion in \citet{waudby-smith_anytime-valid_2024}, and the justifications in \citet{howard_time-uniform_2021}, with a fixed significance level $\alpha$, we have the following CS-width rate:
\begin{equation}
U_t^{\LIL}\left(\pi;\frac{\alpha}{m}\right) - L_t^{\LIL}\left(\pi;\frac{\alpha}{m}\right)
= O\!\left(\frac{\sqrt{\bar V_t(\pi)\left(\log\log \bar V_t(\pi)+\log m\right)}}{t}\right).
\label{eq:lil_rate1}
\end{equation}

\noindent In addition, we will adopt the following assumption.

\begin{assumption}[Linear variance growth]
For each policy $\pi \in \Pi$, the stabilized variance process satisfies
\[
\bar V_t(\pi)=O(t).
\]
\label{lab:assumption_1}
\end{assumption}

\noindent In practice, this assumption says that off-policy evaluation remains statistically well-behaved over time for each policy in the class. This is generally reasonable when the target policy remains sufficiently well-covered by the logging policy, so that large importance weights do not become too frequent or too severe. It can fail when the logging policy assigns probabilities tending to zero to actions that the target policy continues to favor, causing the importance weights---and hence the variance process---to grow too quickly.

\noindent Under Assumption~\ref{lab:assumption_1}, we have $\log\log \bar V_t(\pi)=O(\log\log t)$, and therefore it follows that
\begin{equation}
U_t^{\LIL}\left(\pi;\frac{\alpha}{m}\right) - L_t^{\LIL}\left(\pi;\frac{\alpha}{m}\right)
= O\!\left(\sqrt{\frac{\log\log t+\log m}{t}}\right).
\label{eq:lil_rate2}
\end{equation}
For each fixed policy $\pi \in \Pi$, \eqref{eq:lil_rate2} holds with some policy-dependent constant. Since $\Pi$ is finite, we may take the maximum of these constants over $\Pi$, so the same rate holds uniformly over all $\pi \in \Pi$.

\section{Stopping time definition and rates}

Assume now that $\nu(\pi) \neq \nu(\pi')$ for all $\pi \ne \pi'$, and let $\pi^*$ denote the optimal policy. Define the suboptimality gap for each $\pi \neq \pi^*$ by
\[
\Delta_\pi := \nu(\pi^*) - \nu(\pi),
\]
and define
\[
\Delta_{\min} := \min_{\pi \neq \pi^*}\Delta_\pi > 0.
\]

\begin{theorem}[Stopping time]
Under Assumption~\ref{lab:assumption_1}, fix a significance level $\alpha \in (0,1)$. Define the stopping time
\[
\tau:=\inf\left\{t\ge1:\ \exists \pi'\in\Pi\ \text{s.t.}\ L_t^{\LIL}\left(\pi';\frac{\alpha}{m}\right)\ >\ \max_{\pi \neq \pi'} U_t^{\LIL}\left(\pi;\frac{\alpha}{m}\right)\right\}.
\]
Then, with probability at least $1-\alpha$, $\tau < \infty$, $S_\tau = \{\pi^*\}$, and there exist constants $N>0$, $m_0 \ge 1$, and $\Delta_0 \in (0,1)$ such that for all $m \ge m_0$ and all $0<\Delta_{\min}\le \Delta_0$,
\[
\tau \le N\,\frac{\log m + \log\log(1/\Delta_{\min})}{\Delta_{\min}^2}.
\]
Equivalently,
\[
\tau = O\!\left(\frac{\log m + \log\log(1/\Delta_{\min})}{\Delta_{\min}^2}\right) \qquad \text{as } m\to\infty \text{ and } \Delta_{\min}\downarrow 0.
\]
\label{thm:stopping_time}
\end{theorem}

\subsection{Intermediate lemma}

Before diving into the proof of Theorem~\ref{thm:stopping_time} we will state and prove a small but crucial lemma.

\begin{lemma}[]
Suppose $t \in \mathbb{R}: t \ge C$ for some arbitrary $C > 1$. Additionally, suppose we have constants $A, K \ge e$. Define $L(x):=\log\!\bigl(\log(x)\bigr)$, $X := A + L(K) + L(A) + C$, and $T := 3KXC$. 
It follows that
\[
T > K\bigl(A+L(T)\bigr).
\]
Consequently,
\[
t^\star := \inf\left\{ t \in [C, \infty): t > K\bigl(A+L(t)\bigr) \right\} \le T.
\]
\label{lemma:critical_lemma}
\end{lemma}

\begin{proof}
We begin by establishing that, since $L(K), \ L(A)\ge \log\log(e)=0$, we have $X \ge e + 1$. Similarly, since $X \ge e + 1$ and $K \ge e$, it follows that $T \ge 3(e^2 + e)$. Next, we will show that $L(T)\le X+2$. First, note that by construction $\log \log K := L(K) \le X$ and $C \le X$ which implies that $\log K, \ \log C \le e^X$. Next, since $X \ge e+1$, it follows that $\log X \le e^X$, and $\log 3 \le e^{-(e + 1)}e^X\log 3$, so
\begin{align*}
    \log(T) &=\log3 + \log K + \log X + \log C \\
    &\le e^{-(e + 1)}e^X\log 3 +3e^X \\
    &= (e^{-(e + 1)}\log 3 + 3)e^X.
\end{align*}
It follows that
\begin{align*}
    L(T) &:= \log\log T \\
    &= \log(\log 3 + \log K + \log X + \log C) \\
    &\le \log((e^{-(e + 1)}\log 3 + 3)e^X) \\
    &= \log(e^{-(e + 1)}\log 3 + 3) + X \\
    &< X + 2.
\end{align*}
So, $L(T) \le X + 2$ and as a result $K(A + L(T)) \le K(A + X + 2)$. Finally, since by construction $A \le X$ and $X \ge 3$, we have $2 \le \frac{2X}{3}$ and thus 
\begin{align*}
    A + X + 2 &\le X + X + \frac{2X}{3} \\
    &<3X.
\end{align*}
Combining these pieces, we see that
\begin{align*}
    K(A + L(T)) &\le K(A + X + 2) \\
    &<3KX \\
    &< 3KXC \\
    &= T.
\end{align*}
Since by construction $T > C$, then $t^\star \le T$.
\end{proof}

\subsection{Proof of Theorem~\ref{thm:stopping_time}}

We will now state the proof of Theorem~\ref{thm:stopping_time}.

\begin{proof}
% First, note that $\Delta_{\min} \in (0, 1]$ since $\nu(\pi^*) > \nu(\pi)$ for all $\pi \ne \pi^*$.
Suppose that the simultaneous coverage event holds:
\[
E := \left\{\nu(\pi) \in \left[L_t^\LIL\left(\pi; \frac{\alpha}{m}\right), U_t^\LIL\left(\pi; \frac{\alpha}{m}\right)\right]: \forall t \ge 1, \forall \pi \in \Pi \right\}.
\]
By \eqref{eq:lil_rate2} and the finiteness of $\Pi$, there exist constants $D>0$ and $t_0 > 1$ such that, for every policy $\pi \in \Pi$ and all $t \ge t_0$,
\[
U_t^{\LIL}\left(\pi; \frac{\alpha}{m}\right) - \nu(\pi), \ \nu(\pi) - L_t^{\LIL}\left(\pi; \frac{\alpha}{m}\right)
\le
D\sqrt{\frac{\log\log t+\log m}{t}}.
\]
It immediately follows that, for every suboptimal policy $\pi \neq \pi^*$ and all $t \ge t_0$,
\begin{align*}
L_t^{\LIL}(\pi^*)
-
U_t^{\LIL}(\pi)
&\ge
\nu(\pi^*)
-
D\sqrt{\frac{\log\log t+\log m}{t}}
-
\left(
\nu(\pi)+
D\sqrt{\frac{\log\log t+\log m}{t}}
\right) \\
&=
\Delta_\pi - 2D\sqrt{\frac{\log\log t+\log m}{t}}.
\end{align*}
Therefore, for any $t \ge t_0$,
\[
2D\sqrt{\frac{\log\log t+\log m}{t}} < \Delta_{\min}
\implies
L_t^{\LIL}(\pi^*)\ >\ \max_{\pi \neq \pi^*} U_t^{\LIL}(\pi),
\]
and thus $\tau \le t$. Squaring and rearranging $2D\sqrt{\frac{\log\log t+\log m}{t}} < \Delta_{\min}$ gives us
\[
t>\frac{4D^2}{\Delta_{\min}^2}\bigl(\log m+\log\log t\bigr).
\]
Thus, it suffices to find
\[
t' = \inf\left\{ t \in [t_0, \infty): t > \frac{4D^2}{\Delta_{\min}^2}\bigl(\log m + \log\log t\bigr) \right\}.
\]
It immediately follows that it suffices to find
\[
t^\star = \inf\left\{ t \in [t_0, \infty): t > \max\left\{e, \frac{4D^2}{\Delta_{\min}^2}\right\}\bigl(\max\{e, \log m\} + \log\log t\bigr) \right\},
\]
since
\[
\max\left\{e, \frac{4D^2}{\Delta_{\min}^2}\right\}\bigl(\max\{e, \log m\} + \log\log t\bigr)
\ge
\frac{4D^2}{\Delta_{\min}^2}\bigl(\log m + \log\log t\bigr).
\]
Hence $\tau \le t^\star$. By Lemma~\ref{lemma:critical_lemma}, it also follows that
\begin{align}
\tau
&\le t^\star < T \notag\\
&:= 3t_0\max\left\{e, \frac{4D^2}{\Delta_{\min}^2}\right\}
\left(
\max\{e, \log m\}
+ \log \log \max\left\{e, \frac{4D^2}{\Delta_{\min}^2}\right\}
+ \log \log \max\{e, \log m\}
+ t_0
\right) \notag\\
&\le F\left(
3t_0\frac{4D^2}{\Delta_{\min}^2}
\left(
\log m
+ \log \log \frac{4D^2}{\Delta_{\min}^2}
+ \log \log \max\{e, \log m\}
+ t_0
\right)
\right) \notag\\
&< \infty.
\label{eq:tau_lt_infty}
\end{align}
for a sufficiently large constant $F>0$ depending only on $D$ and $t_0$. Therefore, there exists a constant $M>0$ such that
\[
\tau
\le
M
\frac{
\log m
+ \log \log \left(1/\Delta_{\min}^2\right)
+ \log \log \max\{e, \log m\}
+ 1
}{\Delta_{\min}^2}.
\]
In particular, there exist constants \(m_0 \ge 1\), \(\Delta_0 \in (0,1)\), and \(N>0\) such that for all \(m \ge m_0\) and \(0<\Delta_{\min}\le \Delta_0\),
\[
\tau \le N \frac{\log m + \log\log(1/\Delta_{\min})}{\Delta_{\min}^2}.
\]
That is,
\[
\tau = O\!\left(\frac{\log m + \log\log(1/\Delta_{\min})}{\Delta_{\min}^2}\right)
\qquad \text{as } m\to\infty \text{ and } \Delta_{\min}\downarrow 0.
\]

\noindent Finally, $\tau < \infty$ (\ref{eq:tau_lt_infty}) implies by definition of $S_t$ that $S_\tau = \{\pi^*\}$. By Equation~\ref{eq:simultaneous_coverage}, $\PP(E) \ge 1-\alpha$, and thus all demonstrated results also hold with probability $\ge 1-\alpha$.
\end{proof}

\bibliography{references}

\end{document}